\documentclass[11pt]{article}
\usepackage[T1]{fontenc}
\usepackage{lmodern,microtype}
\usepackage[a4paper,margin=22mm]{geometry}
\usepackage{amsmath,amssymb,amsthm,xcolor}
\usepackage[colorlinks=true,linkcolor=blue!50!black,citecolor=blue!50!black,urlcolor=blue!50!black]{hyperref}
\newtheorem{theorem}{Theorem}
\newtheorem{lemma}[theorem]{Lemma}
\newcommand{\one}{\boldsymbol{1}}
\title{Reliability zeros of maximally connected regular bipartite graphs of large girth}
\author{Research note}
\date{26 November 2026}
\begin{document}
\maketitle
\begin{abstract}
Every point in the closed unit disk is approximated by all-terminal reliability roots of simple regular bipartite graphs with girth and degree tending to infinity and vertex connectivity equal to degree. Every point of $[-1,0]$ can be approximated by real roots. We construct cyclic covers of bipartite expanders in which a matching connects consecutive blocks. A weighted bound on complex component-partition weights permits replacement of exponentially many blocks by single vertices. At any prescribed radius strictly between zero and one, the innermost $d/2$ zeros in a single sequence form an exponentially precise regular polygon, separated from the remaining zeros by a radial gap.
\end{abstract}

For a graph $G=(V,E)$ in which each edge fails independently with probability $z$, let $R_G(z)$ denote the probability of remaining connected. Thus
\[
 R_G(z)=\sum_{A\subseteq E:\,(V,A)\,\mathrm{connected}}
                  (1-z)^{|A|} z^{|E|-|A|},
\]
also defining the complex polynomial. Brown and Colbourn \cite{BC} proved root density in the closed unit disk for multigraphs. The analogous real density conjecture of Brown and McMullin \cite{BM} for simple graphs was settled by Omar \cite{Omar}; Buys \cite{Buys} proved both real and complex density for simple graphs. The clique gadgets in the latter construction give triangles and two-vertex separators. We impose restrictions excluding both features.

\begin{theorem}\label{main}
For each $z_*\in\overline{\mathbb D}$ there exist finite simple bipartite graphs $G_j$ and roots $z_j$ of $R_{G_j}$ with $z_j\to z_*$. We may require that $G_j$ be $d_j$-regular with vertex and edge connectivity $d_j$, and that both $d_j$ and the girth tend to infinity. For $z_*\in[-1,0]$ the roots may all be taken real.
\end{theorem}
The theorem asserts inclusion, not exclusion, for the complex disk. It does not assert density at a single fixed degree, nor do we claim the resulting graphs are expanders.

\section{A maximally connected cyclic cover}
Choose connected, simple, bipartite $d$-regular graphs $X$ of order $n$ with $d\equiv2\pmod4$, $d\to\infty$, $\log n=o(d)$, girth tending to infinity, and second-largest adjacency eigenvalue at most $2\sqrt{d-1}$. Such a sequence follows from Lubotzky, Phillips and Sarnak \cite[Theorems 3.4 and 4.1]{LPS}. Let primes $p\equiv1\pmod4$ tend to infinity; the prime number theorem in this fixed progression supplies a prime $Q\equiv1\pmod4$ between $e^{\sqrt p}$ and $2e^{\sqrt p}$ eventually. Their graph has degree $p+1$, order $Q(Q^2-1)$ or half as large, and girth at least $2\log_p Q$ in this regime. If it is not bipartite, take its connected bipartite double cover, which does not reduce the girth and whose adjacency spectrum consists of both signs of the original spectrum.

Put $k=d/2$, which is odd, and select $k$ edges $a_ib_i$ of a matching with all $a_i$ in the left class. Such a matching exists by Hall's theorem. Remove these edges from $X$ to form $H$; call their endpoints ports. For $m\ge3$ take copies $H_j$ indexed by $j\in\mathbb Z/m\mathbb Z$ and join $a_i$ in $H_j$ to $b_i$ in $H_{j+1}$. Denote the resulting graph by $G$. Its projection onto $X$ is locally bijective on incident edges; hence $G$ is simple, bipartite, $d$-regular and has girth at least that of $X$. Its interblock edges form a matching. For $|U|=u\le n/2$ the centered indicator of $U$ gives $|\partial_XU|\ge(d-2\sqrt{d-1})u(1-u/n)$. Matching deletion costs at most $u$. Thus for sufficiently large $d$,
\begin{equation}\label{expand}
  |\partial_H U|\ge du/4,\qquad \delta(H)=d-1.
\end{equation}

\begin{lemma}The vertex and edge connectivity of $G$ both equal $d$.\end{lemma}
\begin{proof}
Remove $S$ of size at most $d-1$ in one block. Any remaining component $T$ of size $u\le n/2$ satisfies $du/4\le d|S|$, so $u\le4d$. Two vertices have at most one common neighbor since $H$ has no $4$-cycle. Writing $b_y=|N(y)\cap T|$, we obtain $\sum b_y\ge(d-1)u$ and $\sum b_y^2\le du+u(u-1)$. Cauchy--Schwarz bounds the exterior neighborhood of $T$ below by
\[
 f(u)=u\left(\frac{(d-1)^2}{d+u-1}-1\right).
\]
This is concave and $f(2),f(4d)>d-1$ for large $d$. Thus no $2\le u\le4d$ is possible. A singleton component has degree $d-1$ in $H$ and all its neighbors form the deleted set. Two such singletons would have $d-1$ common neighbors. Since $n\gg d$, a disconnected surviving block has one isolated port and one large component.

Now delete fewer than $d$ vertices in the whole graph. If all blocks remain connected, at most one interface vanishes since the $k$ edges per interface and the endpoints of distinct interfaces are disjoint. Otherwise all $d-1$ deletions lie in the block of the isolated port $v$. Its external edge survives. Every deleted vertex has the opposite color from $v$, so the large component contains another port of the same color (there are $k-1$ others), with an intact external edge. The other $m-1$ blocks form a connected path. This proves the lower bound on vertex connectivity and the remaining assertions follow by regularity.
\end{proof}

\section{A weighted component estimate}
For a partition $\pi$ of $V(H)$, write $w_\pi(z)$ for the total weight of configurations with component partition exactly $\pi$. Let $\one$ be the one-part partition and $u(\pi)$ the number of vertices outside a largest part.

\begin{lemma}\label{tail}
For fixed $L\ge1$ and $0<\rho<1$, uniformly on $|z|\le\rho$,
\[
 D_L(z):=\sum_{\pi\ne\one}|w_\pi(z)|L^{u(\pi)}
     =o(\rho^{3d/4}).
\]
\end{lemma}
\begin{proof}
For a partition $\sigma$ let $b(\sigma)$ be the number of edges crossing its parts. All these edges fail precisely when the component partition refines $\sigma$. Consequently by M\"obius inversion
\[
 w_\pi=\sum_{\sigma\le\pi}\mu(\sigma,\pi)z^{b(\sigma)}.
\]
For $u=u(\sigma)>0$ there are at most $n^{2u}$ such $\sigma$: choose the exceptional vertices then partition them. A partition with at most $u+1$ parts has at most $(u+1)^{u+1}\le n^{2u}$ coarsenings. In the usual partition-lattice M\"obius formula, the absolute value is $\prod_B(t_B-1)!$ where $t_B$ is the number of parts merged into $B$. This gives $|\mu(\sigma,\pi)|\le u!\le n^u$. Also $u(\pi)\le u$.

All parts except a largest one have size at most $n/2$. By \eqref{expand} their boundaries give $b(\sigma)\ge du/8$. For $u\le7$ the cut between their union and the largest part has size at least $(d-1)u-u(u-1)=u(d-u)$. These imply $b(\sigma)\ge3d/4+du/32$ for all $u\ge1$ and sufficiently large $d$. Therefore
\[
 D_L(z)\le \rho^{3d/4}\sum_{u\ge1}\left(Ln^5\rho^{d/32}\right)^u
 =o(\rho^{3d/4}),
\]
since $\log n=o(d)$. Also $|w_{\one}-1|\le D_L$ because $\sum_\pi w_\pi=1$.
\end{proof}
\section{Replacing a ring of blocks}
Sum first over the exact partitions $\pi_j$ of each block, then over interblock edge states. This is a polynomial expansion rather than probabilistic conditioning at a complex parameter. If all blocks are connected, the quotient is a cycle with $k$ parallel edges per interface, giving the contribution $w_{\one}^m Q_m(x)$, where $x=z^k$ and
\[
 Q_m(x)=(1-x)^m +mx(1-x)^{m-1}.
\]
We bound every other contribution after grouping configurations that connect the same parts.

Mark a largest part in each block; call an interblock edge exceptional if it touches any unmarked vertex. Since interblock edges form a matching, there are at most $U=\sum_j u(\pi_j)$ exceptional edges and at most $U$ interfaces containing any (\emph{dirty} interfaces). Sum absolute weights for each exceptional edge: at $|z|<1$ the cost per edge is $|z|+|1-z|\le3$. At a dirty interface, the $h$ remaining edges are parallel after contracting marked parts. Their status contributes at most $|1-z^h|+|z^h|\le3$, including $h=0$. At a clean interface the two weights are $1-x$ and $x$. Connectedness permits at most one clean interface to fail completely: two such failures isolate consecutive groups of blocks regardless of other edge states. Thus the conditional interblock sum has absolute value at most
\[
 9^U B,\qquad B=(1+m|x|)e^{m|x|},
\]
by $|1-x|\le1+|x|$ and the choices of zero or one failed clean interface. Putting $D=D_9(z)$ we obtain
\[
 |R_G(z)-w_{\one}^m Q_m(x)|\le
 B\left((|w_{\one}|+D)^m-|w_{\one}|^m\right).
\]
If $mD\to0$ and $m|z|^k$ stays bounded, then $w_{\one}^m\to1$, the displayed right side tends to zero and $Q_m(x)$ stays bounded. Thus $R_G(z)=Q_m(z^k)+o(1)$, uniformly whenever these hypotheses hold uniformly.

\section{The local zero and the physical limit}
Fix $0<r=|z_*|<1$. Take $m-1$ nearest to $r^{-k}$, and a $k$th root $\zeta$ of $-1/(m-1)$ with argument closest to that of $z_*$. Then $\zeta\to z_*$. For negative $z_*$ choose $\zeta$ negative, possible since $k$ is odd. Take $z=\zeta e^{w/k}$ on $|w|\le1/2$. For fixed $r<\rho<r^{2/3}$ we have $|z|\le\rho$ eventually and $m|z|^k$ bounded, while $mD=o(r^{-d/2}\rho^{3d/4})\to0$ by Lemma~\ref{tail}. Since $(m-1)z^k=-e^w$,
\[
  R_G(\zeta e^{w/k})\longrightarrow e^{e^w}(1-e^w)
       \qquad (|w|\le1/2)
\]
uniformly. The limit has exactly one zero there, simple and at zero. Rouch\'e's theorem gives exactly one zero inside for sufficiently large $d$; the corresponding $z$ tends to $z_*$. For negative targets the zero is real by uniqueness and conjugation. Diagonal approximation gives the origin, the unit circle, and both real endpoints. This proves Theorem~\ref{main}. By the general exclusion of Brown and Colbourn \cite[Corollary 3.2]{BC}, the real closure is exactly $[-1,0]\cup\{1\}$.

At a physical failure rate $z=r\in(0,1)$ and $mr^k\to c\in(0,\infty)$, all internal blocks are connected with probability tending to one by Lemma~\ref{tail} and a union bound. The number of fully failed interfaces has law $\operatorname{Binomial}(m,r^k)$. The number of components therefore converges in distribution to $\max(1,Y)$, where $Y\sim\operatorname{Poisson}(c)$, and connectivity has limiting probability $(1+c)e^{-c}$.

\section{An isolated polygon of innermost roots}
In fact a single graph realizes all angular directions at radius $r$, and its other zeros are radially separated.
\begin{theorem}
In the construction of Section 4 put $s=(m-1)^{-1/k}$. There are constants $a,b>0$ depending only on $r$ such that $R_G$ has exactly $k$ zeros in $|z|<s(bd)^{1/k}$ for large $d$. They are simple; there is one at distance $O_r(se^{-ad}/k)$ from each $k$th root of $-1/(m-1)$. Only one of these roots is real and it is negative.
\end{theorem}
\begin{proof}
Take $r<\rho<r^{2/3}$ as above. Then $mD=O_r(e^{-\alpha d})$ for some $\alpha>0$. The bounds of Section 3, $|w_{\one}^m-1|\le e^{mD}-1$ and $|Q_m|\le B$ yield
\[
 |R_G(z)-Q_m(z^k)|\le C_r B e^{-\alpha d}\qquad (|z|\le\rho).
\]
Choose $6b<\alpha$. On $|z|=s(bd)^{1/k}$, eventually inside $|z|\le\rho$, we have $(m-1)|x|=bd$ and $m|x|\le2bd$, where $x=z^k$. Hence
\[
 |Q_m(x)|=|1-x|^{m-1}|1+(m-1)x|\ge e^{-2bd}(bd-1)
\]
by $\log(1-t)\ge-2t$ for $0\le t\le1/2$. This exceeds $C_rBe^{-\alpha d}$ for large $d$. Rouch\'e's theorem counts exactly the $k$ polygon zeros, the other zeros of $Q_m(z^k)$ being on the unit circle.

Around each polygon vertex $\zeta$ write $z=\zeta e^{w/k}$ and let $|w|=e^{-\alpha d/2}$. Uniformly over these contours $B=O(1)$, whereas
\[
 |Q_m(z^k)|=\left|1+\frac{e^w}{m-1}\right|^{m-1}|1-e^w|
 \ge c|w|
\]
for some $c>0$: the first factor tends uniformly to $e$, and the second is at least $|w|/2$. This gives one zero counting multiplicity in each contour by Rouch\'e. The exponential map is injective on each small disk, and their images are disjoint by angular separation $2\pi/k$. Each image has radius $O_r(se^{-\alpha d/2}/k)$ about its center and lies in the preceding larger circle. For odd $k$ the only disk meeting the real line is centered on the negative axis; its unique zero is real by conjugation.
\end{proof}
Thus the next zero has modulus at least $s[1+(\log d+O_r(1))/k]$, whereas the innermost $k$ radii all equal $s$ to exponential precision. No optimality is claimed for the gap.

\paragraph{Girth on the logarithmic scale.}
The witnesses may even satisfy
\[
 \operatorname{girth}(G)\ge \tfrac1{300}\log_{d-1}|V(G)|.
\]
To see this fix $r$, put $\ell=\log(1/r)$, choose $\rho=r^{5/6}$ and in the LPS construction take $Q$ between $e^{p\ell/1000}$ and $2e^{p\ell/1000}$. Here $\log n=(3/1000+o(1))d\ell$, no longer $o(d)$, but the only use of the latter hypothesis for this target was $n^5\rho^{d/32}\to0$. Its logarithm divided by $d\ell$ now tends to $15/1000-5/192<0$, so all estimates used at this $\rho$ remain valid. The order $mn$ has logarithm $(1/2+3/1000+o(1))d\ell$, while the girth is at least $2\log_p Q\ge (2/1000+o(1))d\ell/\log d$. Their ratio on the indicated scale is at least $2/503+o(1)>1/300$ eventually. The constant is independent of $r$, so the endpoint diagonal choices preserve it. This is a fixed positive fraction of the logarithmic scale in the Moore bound, not a claim of an optimal constant.

\small
\begin{thebibliography}{9}
\bibitem{BC} J. I. Brown and C. J. Colbourn, \href{https://doi.org/10.1137/0405047}{Roots of the reliability polynomial}, \emph{SIAM J. Discrete Math.} 5 (1992), 571--585.
\bibitem{BM} J. I. Brown and I. McMullin, \href{https://doi.org/10.1016/j.dam.2026.03.037}{On the real reliability roots of graphs}, \emph{Discrete Appl. Math.} 388 (2026), 174--180.
\bibitem{Buys} P. Buys, \href{https://arxiv.org/abs/2604.07541}{Density of reliability roots of simple graphs in the unit disk}, arXiv:2604.07541 (2026).
\bibitem{LPS} A. Lubotzky, R. Phillips and P. Sarnak, \href{https://doi.org/10.1007/BF02126799}{Ramanujan graphs}, \emph{Combinatorica} 8 (1988), 261--277.
\bibitem{Omar} M. Omar, \href{https://arxiv.org/abs/2604.03530}{Real reliability roots of simple graphs are dense}, arXiv:2604.03530 (2026).
\end{thebibliography}
\end{document}
